\section{Μεγάλα Γλωσσικά Μοντέλα (\en{LLMs}) για Ανάλυση Αρχείων Καταγραφής και Κυβερνοασφάλεια}

\subsection{Εισαγωγή}

Η αυξανόμενη πολυπλοκότητα των σύγχρονων δικτύων, ιδιαίτερα σε περιβάλλοντα \en{5G} και αναδυόμενου \en{6G}, έχει καταστήσει τις παραδοσιακές τεχνικές ανάλυσης αρχείων καταγραφής ανεπαρκείς για αποτελεσματική κυβερνοασφάλεια. Ο τεράστιος όγκος, η ταχύτητα και η μη δομημένη φύση των δεδομένων απαιτούν πιο προηγμένες μεθόδους ανάλυσης που μπορούν να κατανοήσουν συμφραζόμενα και να εντοπίσουν λεπτά πρότυπα. Τα Μεγάλα Γλωσσικά Μοντέλα (\en{LLMs}), με την ικανότητά τους να επεξεργάζονται μη δομημένα δεδομένα κειμένου, έχουν αναδειχθεί ως μια υποσχόμενη λύση.

Τα \en{LLMs} αξιοποιούν βαθιά μάθηση και αρχιτεκτονικές βασισμένες στην προσοχή, καθιστώντας τα κατάλληλα για ανάλυση αρχείων καταγραφής, τα οποία συχνά μοιάζουν με φυσική γλώσσα. Η ικανότητά τους να γενικεύουν σε διαφορετικά περιβάλλοντα επιτρέπει τον εντοπισμό ανωμαλιών με ελάχιστη ανθρώπινη παρέμβαση \cite{zhang2023llmsecurity}.

\subsection{Ανάλυση Αρχείων Καταγραφής με \en{LLMs}}

Η ανάλυση (\en{log parsing}) αποτελεί κρίσιμο στάδιο προεπεξεργασίας, όπου τα ακατέργαστα δεδομένα μετατρέπονται σε δομημένη μορφή. Οι παραδοσιακές μέθοδοι βασίζονται σε κανόνες και δυσκολεύονται να προσαρμοστούν σε νέα μορφότυπα \cite{he2017drain}. Τα \en{LLMs} προσφέρουν πιο ευέλικτη προσέγγιση, αναγνωρίζοντας αυτόματα βασικά στοιχεία όπως τύπους συμβάντων και παραμέτρους. Μπορούν να χειριστούν θόρυβο και ασυνέπειες και να γενικεύσουν σε διαφορετικά συστήματα \cite{liu2023logllm}.

\subsection{Ταξινόμηση Συμβάντων}

Μία βασική εφαρμογή των \en{LLMs} είναι η ταξινόμηση συμβάντων σε κανονικά ή κακόβουλα. Μέσω προσαρμογής (\en{fine-tuning}), τα μοντέλα μαθαίνουν πρότυπα συμπεριφοράς και μπορούν να εντοπίζουν σύνθετες επιθέσεις \cite{nguyen2018deep}. Επιπλέον, μπορούν να εκτελούν πολυκατηγορική ταξινόμηση, αναγνωρίζοντας συγκεκριμένους τύπους επιθέσεων.

\subsection{Πλεονεκτήματα των \en{LLMs} στην Ανάλυση}

Τα \en{LLMs} προσφέρουν σημαντικά πλεονεκτήματα στην ανάλυση δεδομένων και ιδιαίτερα στην επεξεργασία αρχείων καταγραφής, καθώς διαθέτουν ισχυρή ικανότητα κατανόησης συμφραζομένων και σχέσεων μεταξύ στοιχείων σε ακολουθίες δεδομένων. Αυτό τους επιτρέπει να εντοπίζουν σύνθετες επιθέσεις που εκτείνονται σε πολλαπλά συμβάντα και δεν είναι εύκολα ανιχνεύσιμες με παραδοσιακές μεθόδους. Επιπλέον, σε αντίθεση με τις συμβατικές προσεγγίσεις, τα \en{LLMs} μαθαίνουν αυτόματα χρήσιμα χαρακτηριστικά απευθείας από τα δεδομένα, μειώνοντας σημαντικά την ανάγκη για χειροκίνητο \en{feature engineering} και εξειδικευμένη προεπεξεργασία \cite{brown2020language}.

\subsection{Προκλήσεις Χρήσης \en{LLMs}}

Παρά τα σημαντικά τους πλεονεκτήματα, η χρήση των \en{LLMs} συνοδεύεται από ουσιαστικές προκλήσεις. Αρχικά, απαιτούν υψηλούς υπολογιστικούς πόρους τόσο για την εκπαίδευση όσο και για τη λειτουργία τους, γεγονός που δυσχεραίνει την αξιοποίησή τους σε εφαρμογές πραγματικού χρόνου ή σε περιβάλλοντα με περιορισμένες υποδομές. Επιπλέον, η επεξεργασία δεδομένων καταγραφής ενδέχεται να εκθέτει ευαίσθητες πληροφορίες, δημιουργώντας σημαντικά ζητήματα ιδιωτικότητας και κανονιστικής συμμόρφωσης. Για την αντιμετώπιση αυτών των κινδύνων θεωρούνται απαραίτητες τεχνικές όπως η \en{federated learning} και η \en{differential privacy} \cite{xu2021privacy}.